\section{Despliegue y operación de AI Agent}

\subsection{Trade-off entre costo y latencia}

Desplegar un Agent no es solo la selección del modelo, sino también resource orchestration:

\begin{itemize}
    \item \textbf{Arquitectura de inferencia}: chat + tool loop vs. batch prompt, distintas estrategias tienen un impacto enorme en el latency
    \item \textbf{Escalonamiento de recursos}: el path caliente usa low-latency GPU, el path frío usa cheaper CPU
    \item \textbf{Auto-scaling}: escala automáticamente según el request rate, para evitar overprovisioning
    \item \textbf{Multi-model routing}: las tareas de bajo riesgo usan un smaller model, las de alto riesgo usan Sonnet/Opus
\end{itemize}

\begin{importantbox}{No basta con mirar el rendimiento, hay que mirar la end-to-end experience}
El éxito de un Agent no es simplemente la velocidad de inferencia, sino el tiempo total de finalización de la tarea: respuesta del LLM→llamada a la API→posprocesamiento→confirmación del usuario. Un retraso en cualquier paso puede arruinar la experiencia del usuario. El Agent Builder de Vertex AI, mediante un tablero de pipeline, reúne el latency de cada componente de manera concreta en el dashboard.
\end{importantbox}

La lección aquí es que el tamaño del Pod y el prompt size forman dos dimensiones que se presionan mutuamente. Para una token-heavy application, se puede dividir el pipeline en un «first pass de grano grueso» y un «second pass fino», minimizando el costo de model call sin dejar de garantizar el latency.

\subsection{Ciclo de datos y retroalimentación del usuario}

Mejorar continuamente el Agent requiere reutilizar en ciclo el real user data:

\begin{itemize}
    \item \textbf{Structured feedback}: insertar botones de «¿se resolvió el problema?» y «¿cumple con las normas?»
    \item \textbf{Failure tagging}: etiquetar los casos de fallo (hallucination, tool failure, safety violation)
    \item \textbf{Retraining loop}: llevar los failure de alto valor a fine-tuning/LoRA
    \item \textbf{Prompt atlas}: versionar cada modificación de prompt, para facilitar la reversión
\end{itemize}

\begin{knowledgebox}{El Feedback loop da resultados antes que el «reentrenamiento»}
Muchas empresas primero usan el user feedback para ajustar el prompt, y después deciden si reentrenan. Esta estrategia de «primero ajustar la interacción, después hacer un retrain a gran escala» reduce el costo de iteración y permite entregar un Agent confiable más rápido.
\end{knowledgebox}

En Vertex AI, el feedback loop normalmente corre en un `data labeling project` independiente: después de enviar los casos de fallo al labeler para su anotación, se dispara automáticamente el retrain pipeline. Este ciclo automatizado garantiza que el model loop se mantenga fresh en todo momento.

\subsection{Resumen de la sección}

La operación del Agent se enfoca en el costo de latencia, la proporción de recursos y el ciclo de datos; solo después de estandarizar estos elementos el equipo puede llevar las capacidades de laboratorio al entorno de producción.

